% GENERATED FILE. Edit canonical lessons, then run npm run generate:books.
% canonical-source-hash: fnv1a64:1755ce9aa0fb321a
% canonical-lessons: FA-C344-nazr, FA-C344-nazri, FA-C344-halva, FA-C344-shole-zard, FA-C344-golab

\chapter{Vows and sweets}
\label{ch:fa-vows-and-sweets}
\hlchaptermodality{344}

\begin{chapteropening}
\textbf{By the end of this chapter:} \emph{I can talk about a vow, the food given out for it, halva, saffron rice pudding and rosewater.}

Complete the last lesson of chapter 344: I can talk about a vow, the food given out for it, halva, saffron rice pudding and rosewater.
\end{chapteropening}

\section[\texorpdfstring{\fa{نذر} (nazr)}{نذر (nazr)}]{\fa{نذر} (nazr) — a vow (to God)}

\label{lesson:FA-C344-nazr}



\begin{quote}
Before the new one: say the Persian for to cheer, then the Persian for a medal.
\end{quote}

\subsection*{You'll want to know: \fa{نذر}}
\textbf{\fa{نذر}} (\emph{nazr}) — \textquotedblleft{}a vow (to God)\textquotedblright{}.

\begin{quote}
\fa{او} \fa{نذر} \fa{کرد} \fa{که} \fa{اگر} \fa{حالش} \fa{خوب} \fa{شود} \fa{به} \fa{فقیرها} \fa{کمک} \fa{کند}.
\end{quote}

\begin{quote}
\emph{U nazr kard ke agar hâlash khub shavad be faqir-hâ komak konad.}
\end{quote}

\begin{quote}
She vowed that if she got well she would help the poor.
\end{quote}

Written with \fa{ذ}. \emph{Nazr kardan} \textquotedblleft{}to make a vow\textquotedblright{}; it is kept by cooking or giving.

\begin{cousinweb}
\emph{nazr} is borrowed from Arabic \emph{nadhr}, \textquotedblleft{}a vow\textquotedblright{} (EtymDB); Persian says the \fa{ذ} as \emph{z}.
\end{cousinweb}

\begin{grammarlens}[title={What you've built}]
One more word for telling, explaining and describing at B1.
\end{grammarlens}

\subsection*{Your turn}
\begin{itemize}
\raggedright
  \item \emph{Say it:} \emph{nazr}
  \item \emph{Say it:} \emph{nazr}, once more
  \item \emph{Say it:} a sentence of your own with \emph{nazr}
  \item \emph{Recall:} say the Persian for to cheer, then the Persian for a medal, then say \emph{nazr} again
\end{itemize}

\begin{tcolorbox}[breakable,colback=teal!4,colframe=teal!35!black,title={Before you move on}]
What does \fa{نذر} mean? (\textquotedblleft{}A vow (to God)\textquotedblright{}.) Say it once more.
\end{tcolorbox}
\section[\texorpdfstring{\fa{نذری} (nazri)}{نذری (nazri)}]{\fa{نذری} (nazri) — votive food, food given out for a vow}

\label{lesson:FA-C344-nazri}



\begin{quote}
Before the new one: say the Persian for a medal, then the Persian for a vow.
\end{quote}

\subsection*{You'll want to know: \fa{نذری}}
\textbf{\fa{نذری}} (\emph{nazri}) — \textquotedblleft{}votive food, food given out for a vow\textquotedblright{}.

\begin{quote}
\fa{همسایه} \fa{ها} \fa{برای} \fa{همه} \fa{نذری} \fa{پخش} \fa{کردند}.
\end{quote}

\begin{quote}
\emph{Hamsâye-hâ barâ-ye hame nazri pakhsh kardand.}
\end{quote}

\begin{quote}
The neighbours handed out votive food to everyone.
\end{quote}

\emph{nazr} + \emph{-i}: food cooked to keep a vow and shared with neighbours and strangers.

\begin{cousinweb}
\emph{nazri} adds the Persian suffix \emph{-i} to \emph{nazr} \textquotedblleft{}a vow\textquotedblright{}, which is borrowed from Arabic \emph{nadhr} (EtymDB).
\end{cousinweb}

\begin{grammarlens}[title={What you've built}]
One more word for telling, explaining and describing at B1.
\end{grammarlens}

\subsection*{Your turn}
\begin{itemize}
\raggedright
  \item \emph{Say it:} \emph{nazri}
  \item \emph{Say it:} \emph{nazri}, once more
  \item \emph{Say it:} a sentence of your own with \emph{nazri}
  \item \emph{Recall:} say the Persian for a medal, then the Persian for a vow, then say \emph{nazri} again
\end{itemize}

\begin{tcolorbox}[breakable,colback=teal!4,colframe=teal!35!black,title={Before you move on}]
What does \fa{نذری} mean? (\textquotedblleft{}Votive food, food given out for a vow\textquotedblright{}.) Say it once more.
\end{tcolorbox}
\section[\texorpdfstring{\fa{حلوا} (halvâ)}{حلوا (halvâ)}]{\fa{حلوا} (halvâ) — halva}

\label{lesson:FA-C344-halva}



\begin{quote}
Before the new one: say the Persian for a vow, then the Persian for votive food.
\end{quote}

\subsection*{You'll want to know: \fa{حلوا}}
\textbf{\fa{حلوا}} (\emph{halvâ}) — \textquotedblleft{}halva\textquotedblright{}.

\begin{quote}
\fa{در} \fa{مراسم} \fa{ختم} \fa{حلوا} \fa{پخش} \fa{کردند}.
\end{quote}

\begin{quote}
\emph{Dar marâsem-e khatm halvâ pakhsh kardand.}
\end{quote}

\begin{quote}
At the memorial service they handed out halva.
\end{quote}

Persian halva is flour, butter, sugar and saffron, served at mourning ceremonies.

\begin{cousinweb}
\emph{halvâ} is borrowed from Arabic \emph{ḥalwā} (EtymDB), which EtymDB derives from \emph{ḥulw} \textquotedblleft{}sweet\textquotedblright{}.
\end{cousinweb}

\begin{grammarlens}[title={What you've built}]
One more word for telling, explaining and describing at B1.
\end{grammarlens}

\subsection*{Your turn}
\begin{itemize}
\raggedright
  \item \emph{Say it:} \emph{halvâ}
  \item \emph{Say it:} \emph{halvâ}, once more
  \item \emph{Say it:} a sentence of your own with \emph{halvâ}
  \item \emph{Recall:} say the Persian for a vow, then the Persian for votive food, then say \emph{halvâ} again
\end{itemize}

\begin{tcolorbox}[breakable,colback=teal!4,colframe=teal!35!black,title={Before you move on}]
What does \fa{حلوا} mean? (\textquotedblleft{}Halva\textquotedblright{}.) Say it once more.
\end{tcolorbox}
\section[\texorpdfstring{\fa{شله} \fa{زرد} (shole zard)}{شله زرد (shole zard)}]{\fa{شله} \fa{زرد} (shole zard) — saffron rice pudding}

\label{lesson:FA-C344-shole-zard}



\begin{quote}
Before the new one: say the Persian for votive food, then the Persian for halva.
\end{quote}

\subsection*{You'll want to know: \fa{شله} \fa{زرد}}
\textbf{\fa{شله} \fa{زرد}} (\emph{shole zard}) — \textquotedblleft{}saffron rice pudding\textquotedblright{}.

\begin{quote}
\fa{مادرم} \fa{برای} \fa{نذری} \fa{شله} \fa{زرد} \fa{پخت}.
\end{quote}

\begin{quote}
\emph{Mâdaram barâ-ye nazri shole zard pokht.}
\end{quote}

\begin{quote}
My mother cooked saffron rice pudding as votive food.
\end{quote}

Rice, sugar, saffron and rosewater, decorated with cinnamon. \emph{shole} \textquotedblleft{}a soft rice dish\textquotedblright{} + \emph{zard} \textquotedblleft{}yellow\textquotedblright{}.

\begin{cousinweb}
\emph{shole zard} joins \emph{shole} \textquotedblleft{}a soft rice dish\textquotedblright{} and \emph{zard} \textquotedblleft{}yellow, golden\textquotedblright{}, a native Persian word that EtymDB lists without a parent.
\end{cousinweb}

\begin{grammarlens}[title={What you've built}]
One more word for telling, explaining and describing at B1.
\end{grammarlens}

\subsection*{Your turn}
\begin{itemize}
\raggedright
  \item \emph{Say it:} \emph{shole zard}
  \item \emph{Say it:} \emph{shole zard}, once more
  \item \emph{Say it:} a sentence of your own with \emph{shole zard}
  \item \emph{Recall:} say the Persian for votive food, then the Persian for halva, then say \emph{shole zard} again
\end{itemize}

\begin{tcolorbox}[breakable,colback=teal!4,colframe=teal!35!black,title={Before you move on}]
What does \fa{شله} \fa{زرد} mean? (\textquotedblleft{}Saffron rice pudding\textquotedblright{}.) Say it once more.
\end{tcolorbox}
\section[\texorpdfstring{\fa{گلاب} (golâb)}{گلاب (golâb)}]{\fa{گلاب} (golâb) — rosewater}

\label{lesson:FA-C344-golab}



\begin{quote}
Before the new one: say the Persian for halva, then the Persian for saffron rice pudding.
\end{quote}

\subsection*{You'll want to know: \fa{گلاب}}
\textbf{\fa{گلاب}} (\emph{golâb}) — \textquotedblleft{}rosewater\textquotedblright{}.

\begin{quote}
\fa{بوی} \fa{گلاب} \fa{در} \fa{همه} \fa{خانه} \fa{بود}.
\end{quote}

\begin{quote}
\emph{Bu-ye golâb dar hame-ye khâne bud.}
\end{quote}

\begin{quote}
The scent of rosewater filled the whole house.
\end{quote}

\emph{gol} \textquotedblleft{}rose, flower\textquotedblright{} + \emph{âb} \textquotedblleft{}water\textquotedblright{}.

\begin{cousinweb}
\emph{golâb} is a compound of \emph{gol} \textquotedblleft{}rose\textquotedblright{} and \emph{âb} \textquotedblleft{}water\textquotedblright{}, both native Persian words (EtymDB lists the compound).
\end{cousinweb}

\begin{grammarlens}[title={What you've built}]
One more word for telling, explaining and describing at B1.
\end{grammarlens}

\subsection*{Your turn}
\begin{itemize}
\raggedright
  \item \emph{Say it:} \emph{golâb}
  \item \emph{Say it:} \emph{golâb}, once more
  \item \emph{Say it:} a sentence of your own with \emph{golâb}
  \item \emph{Recall:} say the Persian for halva, then the Persian for saffron rice pudding, then say \emph{golâb} again
\end{itemize}

\begin{tcolorbox}[breakable,colback=teal!4,colframe=teal!35!black,title={Before you move on}]
What does \fa{گلاب} mean? (\textquotedblleft{}Rosewater\textquotedblright{}.) Say it once more.
\end{tcolorbox}
